\editionPage[\pdfbookmark[0]{Indications des défis}{bm-indices-debut-parent}]{indices-debut}{Indications progressives}{Défis 1 à 3}{Ne lire que l’aide nécessaire.}
\begin{blocEdition}{indices-debut}
\editionEtape{D01 · Cent termes, sans cent calculs}\textbf{Indice 1.} Regrouper les termes par paires consécutives.\par\medskip\textbf{Indice 2.} Chaque paire $-(2k-1)+2k$ vaut 1.\par
\editionEtape{D02 · Retrouver tout le budget}\textbf{Indice 1.} Exprimer chaque reste comme fraction du budget initial.\par\medskip\textbf{Indice 2.} Le reste final représente $(2/3)(3/4)(2/5)$ du total.\par
\editionEtape{D03 · Jusqu’où peut-on doubler ?}\textbf{Indice 1.} Après $n$ étapes, combien de facteurs 2 a-t-on introduits ?\par\medskip\textbf{Indice 2.} Calculer $2^{10}$, puis comparer $2^{19}$ et $2^{20}$ à un million.\par

\end{blocEdition}
\editionPage{indices-4}{Indications progressives}{Défis 4 à 6}{Ne lire que l’aide nécessaire.}
\begin{blocEdition}{indices-4}
\editionEtape{D04 · Quatre carreaux pour un grand carré}\textbf{Indice 1.} Le côté du grand carré contient deux côtés de carreaux.\par\medskip\textbf{Indice 2.} Son aire est la somme des quatre aires ; sa longueur est aussi $\sqrt{72}$.\par
\editionEtape{D05 · Comparer sans choisir de valeur}\textbf{Indice 1.} Développer les deux expressions d’aire.\par\medskip\textbf{Indice 2.} Calculer l’aire de A moins celle de B.\par
\editionEtape{D06 · Compter les pièces}\textbf{Indice 1.} Appeler $x$ le nombre de pièces de 2 euros.\par\medskip\textbf{Indice 2.} Il reste $42-x$ pièces de 1 euro.\par

\end{blocEdition}
\editionPage{indices-7}{Indications progressives}{Défis 7 à 9}{Ne lire que l’aide nécessaire.}
\begin{blocEdition}{indices-7}
\editionEtape{D07 · Annuler une remise}\textbf{Indice 1.} Après la baisse, il reste 80\,\% du prix initial.\par\medskip\textbf{Indice 2.} La somme à rajouter représente 20 parts sur les 80 parts du nouveau prix.\par
\editionEtape{D08 · Un remplissage avec fuite}\textbf{Indice 1.} Distinguer les deux périodes.\par\medskip\textbf{Indice 2.} Après 20 minutes, le volume a augmenté de $20(18-3)$ litres.\par
\editionEtape{D09 · Un découpage variable}\textbf{Indice 1.} L’aire retirée est $4x^2$.\par\medskip\textbf{Indice 2.} Si l’aire restante est 116, chaque carré retiré a une aire de 16.\par

\end{blocEdition}
\editionPage{indices-10}{Indications progressives}{Défis 10 à 12}{Ne lire que l’aide nécessaire.}
\begin{blocEdition}{indices-10}
\editionEtape{D10 · Même moyenne, autre médiane}\textbf{Indice 1.} Les deux sommes valent 60.\par\medskip\textbf{Indice 2.} Avec six valeurs, observer les rangs 3 et 4.\par
\editionEtape{D11 · Une cible et une comparaison exacte}\textbf{Indice 1.} Comparer l’aire du disque à celle du carré.\par\medskip\textbf{Indice 2.} Le quotient se simplifie en $\pi/16$.\par
\editionEtape{D12 · Une droite, trois rôles}\textbf{Indice 1.} Comparer ABM et ACM avec trois côtés.\par\medskip\textbf{Indice 2.} Les deux angles en M sont égaux et supplémentaires.\par

\end{blocEdition}
\editionPage{indices-13}{Indications progressives}{Défis 13 à 15}{Ne lire que l’aide nécessaire.}
\begin{blocEdition}{indices-13}
\editionEtape{D13 · Le câble et ses attaches}\textbf{Indice 1.} Calculer d’abord la distance entre les points d’attache.\par\medskip\textbf{Indice 2.} Deux extrémités demandent deux réserves de 0,75 m.\par
\editionEtape{D14 · Une bande entre triangles}\textbf{Indice 1.} Le rapport petit sur grand vaut $1/4$.\par\medskip\textbf{Indice 2.} Calculer BC par Pythagore et les deux aires séparément.\par
\editionEtape{D15 · Choisir la bonne inclinaison}\textbf{Indice 1.} Écrire le cosinus avec chaque longueur horizontale.\par\medskip\textbf{Indice 2.} Utiliser le cosinus inverse en degrés pour les trois rapports.\par

\end{blocEdition}
\editionPage{indices-16}{Indications progressives}{Défis 16 à 18}{Ne lire que l’aide nécessaire.}
\begin{blocEdition}{indices-16}
\editionEtape{D16 · Deux glissements, une construction}\textbf{Indice 1.} Additionner séparément les déplacements horizontaux et verticaux.\par\medskip\textbf{Indice 2.} Le déplacement total est de 3 à droite et 3 en haut.\par
\editionEtape{D17 · Un volume sous deux formes}\textbf{Indice 1.} Les deux volumes diffèrent par le facteur $1/3$.\par\medskip\textbf{Indice 2.} Réduire les longueurs de moitié divise le volume par 8.\par
\editionEtape{D18 · Une somme programmée et prouvée}\textbf{Indice 1.} n vaut le prochain entier à ajouter.\par\medskip\textbf{Indice 2.} Regrouper $1+10$, $2+9$, etc.\par

\end{blocEdition}
